\section{Selection mathematics: Price and landscapes}
\label{sec:selection}

\subsection{Price equation (structural)}
Price's covariance formulation decomposes change in average trait value into selection and transmission components~\cite{price1970selection,frank1995george}.
For measurable trait $z_i$ and fitness $w_i$:
\begin{equation}
  \Delta \bar{z}
  = \frac{\mathrm{Cov}(w_i,z_i)}{\bar{w}}
  + \frac{\mathbb{E}[w_i\,\Delta z_i]}{\bar{w}}.
\end{equation}
The first term captures directional selection; the second captures mutation and inheritance biases.

\subsection{Fitness landscapes}
A genotype $g\in\mathcal{G}$ maps to fitness $F(g)$.
Epistasis (non-additive interactions) creates multimodal landscapes where greedy hill climbing fails~\cite{mitchell1998introduction}.
This motivates crossover (recombination), niching, and quality diversity~\cite{deb2002fast,mouret2015illuminating}.

\subsection{Fisher's fundamental theorem: caution}
Fisher's theorem has generated extensive clarification because informal statements often misrepresent scope~\cite{fisher1930genetical}.
This thesis treats Fisher primarily as historical context and relies on formal bandit or convergence statements when precise guarantees are needed~\cite{lai1985asymptotically,auer2002finite}.
